\chapter{Arquitetura do sistema multiagente}
\label{cap:arquitetura}

% =====================================================================
% PAPEL DESTE CAPÍTULO (registro de projeto):
%   Apresenta a arquitetura em nível conceitual: papéis, contratos de
%   interface e fluxo de execução. NÃO repete a justificativa da
%   combinação (casa canônica: 2.5) nem a fundamentação da rubrica
%   (casa canônica: 3.4). NÃO desce a detalhes de código (Cap. 5).
% =====================================================================

Os capítulos anteriores estabeleceram o problema (Capítulo~\ref{cap:contexto}), a hipótese
arquitetural que responde a ele (Seção~\ref{sec:hipotese}) e os referenciais didáticos que
fundamentam os critérios de qualidade (Seção~\ref{sec:criterios}). Este capítulo apresenta o
projeto do sistema que materializa essa hipótese: a especificação de entrada fornecida pelo
professor, os quatro agentes com seus contratos de interface, o fluxo de execução completo e as
decisões de projeto --- incluindo as alternativas consideradas e descartadas. A implementação
concreta, com as escolhas de tecnologia e trechos de código, é objeto do
Capítulo~\ref{cap:implementacao}.

\section{Visão geral}
\label{sec:visao-geral}

O sistema é composto por quatro agentes especializados coordenados em um fluxo cíclico de
geração e revisão, conforme a Figura~\ref{fig:componentes}. O professor fornece uma
\textbf{especificação} --- os parâmetros didáticos da questão desejada --- e recebe, ao final,
ou uma questão aprovada acompanhada dos veredictos de todos os avaliadores, ou a informação de
que a especificação foi descartada após esgotar as tentativas de revisão.

\begin{figure}[htbp]
\centering
\begin{tabular}{c}
\begin{minipage}{0.92\textwidth}
\centering\ttfamily\small
[Especificação do professor]\\[2pt]
$\downarrow$\\[2pt]
\fbox{Orquestrador} $\rightarrow$ \fbox{Gerador (LLM)} $\rightarrow$ questão estruturada\\[2pt]
$\downarrow$\\[2pt]
\fbox{Verificador Simbólico (SymPy)} \quad e \quad \fbox{Crítico Didático (LLM + rubrica)}\\[2pt]
$\downarrow$\\[2pt]
aprovada $\rightarrow$ banco curado \qquad
rejeitada $\rightarrow$ Gerador (com \emph{feedback}) ou descarte
\end{minipage}
\end{tabular}
\caption{Diagrama de componentes do sistema. Elaborado pela autora.}
\label{fig:componentes}
\end{figure}

Três princípios de projeto, derivados diretamente da fundamentação técnica, governam a
arquitetura. O primeiro é a \textbf{separação de responsabilidades}: cada agente executa uma
função única e bem delimitada, e nenhum agente acumula os papéis de gerar e avaliar
(Seção~\ref{sub:sma}). O segundo é a \textbf{verificação independente por mecanismo distinto}: a
correção matemática é julgada por um sistema de computação algébrica, não por um segundo LLM,
de modo que os erros dos dois componentes tenham naturezas estruturalmente diferentes
(Seção~\ref{sec:hipotese}). O terceiro é a \textbf{rastreabilidade}: toda iteração do ciclo ---
incluindo as tentativas rejeitadas --- é registrada e apresentada ao professor, que recebe não
apenas o produto final, mas a história da sua produção.

\section{Especificação de entrada}
\label{sec:especificacao}

A especificação é o contrato entre o professor e o sistema. Ela é composta por seis parâmetros
obrigatórios e dois opcionais, todos validados antes do início do ciclo:

\begin{itemize}
\item \textbf{Habilidade BNCC} --- o código da habilidade (por exemplo, \texttt{EM13MAT302}),
  validado contra um catálogo local derivado da BNCC; códigos desconhecidos são rejeitados na
  entrada, antes de qualquer chamada ao LLM. É o \emph{primeiro} parâmetro, e não por acaso: a
  habilidade é a unidade normativa do planejamento do professor e é ela que delimita os temas
  disponíveis (Seção~\ref{sub:bncc-habilidades}).
\item \textbf{Tema ou temas} --- o tópico matemático, escolhido \emph{dentro} da habilidade:
  função afim, função quadrática, função exponencial, progressão aritmética ou progressão
  geométrica. Uma habilidade pode cobrir mais de um tema, e o catálogo registra como eles se
  combinam. Em \texttt{EM13MAT302} a enumeração é do repertório coberto, e combinar os dois
  temas num mesmo problema é opção do professor; já em \texttt{EM13MAT507}, cujo texto é
  ``identificar e associar sequências numéricas (PA) a funções afins de domínios discretos'', a
  articulação entre os temas \emph{é} a habilidade --- pedir apenas a progressão a
  descaracteriza, e o sistema recusa a especificação. Com mais de um tema, a questão pedida é
  uma só, articulando-os, e não dois itens justapostos.
\item \textbf{Nível cognitivo} --- um dos seis processos da Taxonomia de Bloom Revisada
  (Seção~\ref{sub:bloom}). O catálogo registra os níveis compatíveis com os verbos de cada
  habilidade; pedir um nível fora deles é permitido --- o professor pode ter suas razões ---, mas
  a interface avisa e a divergência fica registrada no ciclo, como dado da avaliação empírica.
\item \textbf{Dificuldade} --- fácil, média ou difícil, com descritores operacionais que
  acompanham o \emph{prompt} (fácil: aplicação direta; média: articulação de dois ou mais
  passos; difícil: estratégia não evidente ou análise de casos).
\item \textbf{Natureza} --- teórica (manipulação e propriedades, sem contexto externo) ou
  aplicada (situação-problema com contexto realista).
\item \textbf{Formato} --- discursiva ou múltipla escolha com quatro alternativas.
\item \textbf{Contexto temático} (opcional) --- um domínio para ambientar a questão aplicada
  (esportes, finanças pessoais etc.).
\item \textbf{Restrições} (opcionais) --- condições adicionais em linguagem natural (por
  exemplo, ``apenas coeficientes inteiros'').
\end{itemize}

A validação antecipada do código BNCC merece destaque: ela impede, por construção, que o
sistema produza questões associadas a habilidades inexistentes --- um erro que um LLM cometeria
com naturalidade e que seria invisível ao professor que não tem o documento da BNCC de memória.

A habilidade, porém, não entra no sistema apenas como rótulo a ser validado. O catálogo associa
a cada habilidade um conjunto de \textbf{exigências}: o que a questão precisa efetivamente
exibir para realizá-la. Para \texttt{EM13MAT402} --- ``converter representações algébricas de
funções polinomiais de 2º grau para representações geométricas no plano cartesiano'' --- exige-se
que a questão faça o estudante transitar entre a forma algébrica e o gráfico; pedir apenas as
raízes não realiza a habilidade, ainda que o conteúdo seja o correto e a matemática esteja certa.
Essas exigências são repassadas ao Gerador como requisito de elaboração e ao Crítico como
âncora do critério de alinhamento curricular. Sem elas, o alinhamento à BNCC dependeria de um
juízo em prosa do LLM sobre se a questão ``é sobre o assunto'' --- e uma questão correta,
clara e sobre o conteúdo certo poderia ser aprovada sem cumprir a habilidade declarada, um
erro silencioso de natureza curricular, análogo ao erro silencioso matemático discutido na
Seção~\ref{sec:erro-silencioso}.

\section{Agente Gerador}
\label{sec:ag-gerador}

O Gerador é o único agente criativo do sistema. Recebe a especificação (e, nas iterações de
revisão, o \emph{feedback} estruturado do Orquestrador) e produz a questão completa como objeto
estruturado, nunca como texto livre. Sua saída contém: o enunciado; a resolução passo a passo;
o gabarito em linguagem natural; as alternativas, quando múltipla escolha, cada distrator
anotado com o erro sistemático de estudante que ele representa; e o campo mais importante do
ponto de vista arquitetural, a \textbf{formalização verificável}.

A formalização verificável é a ponte entre a prosa da questão e o verificador simbólico. Trata-se
de um objeto com quatro campos --- o \emph{tipo} de verificação (equação, propriedade de função
ou progressão), a \emph{expressão} em sintaxe SymPy, a \emph{resposta esperada} (o gabarito, na
mesma sintaxe) e \emph{parâmetros} auxiliares --- que o Gerador preenche no momento da criação da
questão. Essa decisão resolve o segundo limite da verificação simbólica apontado na
Seção~\ref{sub:cas-limites}: o Verificador nunca precisa interpretar o enunciado em linguagem
natural, pois quem traduz a prosa para a forma simbólica é o próprio agente que a escreveu ---
e a tradução é, ela mesma, parte do objeto auditável que o professor pode inspecionar. Quando a
questão não admite formalização nos tipos suportados, o campo é nulo e a questão segue o caminho
do veredicto \emph{não-verificável} (Seção~\ref{sec:ag-verificador}).

O \emph{prompt} do Gerador é um artefato versionado da pesquisa (reproduzido no Apêndice A) e
possui três blocos: a instrução de papel (elaborador experiente de questões para o EM), as
regras obrigatórias --- incluindo a exigência de saída exclusivamente em JSON e a proibição de
sinalizar o procedimento de resolução no enunciado, em atenção ao efeito Topaze
(Seção~\ref{sub:tsd}) --- e o esquema do objeto de saída.

\section{Agente Verificador Simbólico}
\label{sec:ag-verificador}

O Verificador resolve a questão de forma independente e compara o resultado com o gabarito
proposto. É o único agente sem LLM: opera exclusivamente sobre a formalização verificável, por
meio do SymPy, e por isso seu julgamento é determinístico e reprodutível
(Seção~\ref{sec:cas}).

A verificação é organizada por \textbf{estratégias por tipo de questão}:

\begin{itemize}
\item \textbf{Equações} --- o conjunto-solução é calculado simbolicamente e comparado com o
  gabarito por igualdade exata (a diferença entre cada par de soluções deve simplificar a zero).
  A comparação detecta tanto soluções incorretas quanto \emph{gabaritos incompletos} --- uma
  raiz esquecida, erro frequente de LLMs.
\item \textbf{Propriedades de funções} --- zeros, vértice, valor em um ponto e extremos, cada um
  calculado pela via simbólica adequada (resolução, derivada, substituição).
\item \textbf{Progressões} --- termo geral e soma de PA e PG, calculados pelas fórmulas fechadas,
  incluindo o caso degenerado da PG de razão unitária.
\item \textbf{Verificação numérica complementar} --- quando a simplificação simbólica não é
  conclusiva, as expressões são avaliadas em pontos aleatórios com semente fixa; a concordância
  gera o veredicto \emph{aprovado com ressalva numérica} (Seção~\ref{sub:cas-limites}).
\end{itemize}

A saída é um veredicto de quatro valores --- \emph{aprovado}, \emph{rejeitado},
\emph{não-verificável} ou \emph{aprovado com ressalva numérica} --- acompanhado de justificativa
e, quando houver, do resultado calculado independentemente. Esse resultado calculado integra o
\emph{feedback} devolvido ao Gerador em caso de rejeição: o agente que errou recebe não apenas a
informação de que errou, mas o valor correto contra o qual revisar seu trabalho.

Um princípio de robustez completa o projeto deste agente: formalizações malformadas (sintaxe
inválida produzida pelo LLM, tipo desconhecido) nunca interrompem o ciclo com erro --- degradam
para o veredicto \emph{não-verificável}, e a decisão sobre o que fazer com a questão permanece
onde deve estar, no Orquestrador.

\section{Agente Crítico Didático}
\label{sec:ag-critico}

O Crítico avalia a qualidade pedagógica da questão com base na rubrica consolidada na
Seção~\ref{sec:criterios}. Como a correção matemática já foi julgada pelo Verificador, o
Crítico avalia os cinco critérios restantes --- clareza do enunciado, adequação ao nível
cognitivo, alinhamento à habilidade BNCC, plausibilidade dos distratores e originalidade ---
atribuindo a cada um nota de 1 a 5 e comentário objetivo.

O critério de alinhamento à BNCC recebe tratamento distinto dos demais: em vez de pedir um juízo
global sobre o assunto da questão, o \emph{prompt} apresenta ao Crítico as exigências da
habilidade declarada (Seção~\ref{sec:especificacao}) e o instrui a julgá-las uma a uma. Abordar o
conteúdo certo não basta: se a questão não faz o que a habilidade pede, a nota é baixa ainda que
o enunciado esteja claro e a matemática correta.

A rubrica é traduzida em \emph{prompt} (Apêndice A) com uma instrução de decisão explícita: a
questão é aprovada somente se nenhum critério receber nota inferior a 3. Em caso de reprovação,
o Crítico é obrigado a preencher um campo de \textbf{sugestões de revisão} com instruções
acionáveis --- o que mudar, e não apenas o que está errado --- que o Orquestrador repassa ao
Gerador. O parecer completo (notas, comentários, decisão e sugestões) é um objeto estruturado
que acompanha a questão até o professor.

\section{Agente Orquestrador}
\label{sec:ag-orquestrador}

O Orquestrador coordena o ciclo e concentra toda a política de decisão do sistema. Ele não
gera, não verifica e não avalia: decide. Sua política, deliberadamente simples e auditável, é:

\begin{enumerate}
\item \textbf{Rejeição do Verificador} --- a questão volta ao Gerador com o veredicto, a
  justificativa e o resultado calculado; \emph{não passa pelo Crítico}, pois não há sentido em
  avaliar didaticamente uma questão matematicamente incorreta.
\item \textbf{Veredicto não-verificável} --- a questão \emph{segue} ao Crítico. A
  impossibilidade de verificação automática não é defeito da questão; a aprovação final carrega
  a marca de revisão prioritária pelo professor.
\item \textbf{Reprovação do Crítico} --- a questão volta ao Gerador com as sugestões de revisão.
\item \textbf{Aprovação dos dois avaliadores} --- a questão é aprovada e entra no banco curado.
\item \textbf{Critério de parada} --- três iterações sem aprovação implicam descarte. O limite
  evita ciclos improdutivos (e o custo correspondente) quando a especificação é inviável ou o
  Gerador não converge.
\end{enumerate}

A \textbf{memória de execução} completa o papel do agente: cada iteração --- questão gerada,
veredicto do Verificador, parecer do Crítico e \emph{feedback} enviado --- é registrada em um
objeto de resultado que acompanha o ciclo inteiro. Esse registro cumpre dupla função:
rastreabilidade para o professor (que pode inspecionar por que uma questão precisou de duas
tentativas) e material de análise para a avaliação empírica do sistema.

\section{Fluxo de execução completo}
\label{sec:fluxo}

O percurso típico de uma especificação pelo sistema pode ser descrito em cinco tempos:

\begin{enumerate}
\item O professor submete a especificação; a validação de entrada rejeita imediatamente códigos
  BNCC desconhecidos.
\item O Orquestrador aciona o Gerador, que produz a questão estruturada com sua formalização
  verificável.
\item O Verificador resolve a formalização e emite seu veredicto; se rejeitar, o ciclo retorna
  ao passo 2 com \emph{feedback}, sem consultar o Crítico.
\item O Crítico avalia a questão pela rubrica; se reprovar, o ciclo retorna ao passo 2 com as
  sugestões de revisão.
\item Aprovada por ambos --- ou esgotadas as três iterações --- o Orquestrador registra o
  resultado completo, e a questão aprovada é persistida no banco curado com seus metadados e
  histórico.
\end{enumerate}

Dois exemplos ilustram o valor do ciclo. No primeiro, o Gerador produz uma questão de progressão
geométrica cujo gabarito omite o primeiro termo ($3^5 = 243$ em vez de $2 \cdot 3^5 = 486$) ---
um erro silencioso clássico, embutido em uma resolução fluente. O Verificador calcula o termo de
forma independente, rejeita o gabarito e devolve o valor correto; a segunda iteração é aprovada.
No segundo exemplo, uma questão matematicamente impecável é reprovada pelo Crítico por
inconsistência de nível: especificada como \emph{analisar}, exigia apenas a aplicação de um
algoritmo anunciado no próprio enunciado; a revisão eleva a demanda cognitiva. Os dois casos ---
detalhados, com os objetos completos, no Capítulo~\ref{cap:implementacao} --- correspondem
exatamente às duas dimensões de qualidade que a hipótese arquitetural (Seção~\ref{sec:hipotese})
se propõe a garantir.

\section{Decisões de projeto e alternativas descartadas}
\label{sec:decisoes}

Quatro decisões estruturais merecem justificativa explícita.

\textbf{Por que quatro agentes, e não três ou cinco.} A divisão reflete as quatro funções
irredutíveis do problema: criar, verificar matemática, avaliar didática e decidir. Fundir
verificador e crítico em um único avaliador destruiria a propriedade central da arquitetura ---
a independência de mecanismos (Seção~\ref{sec:hipotese}); adicionar um quinto agente (por
exemplo, um revisor de linguagem) aumentaria custo e latência sem função que os existentes não
cubram.

\textbf{Por que verificação independente, e não auto-verificação.} Solicitar ao próprio LLM que
confira seu trabalho é sujeito aos mesmos vieses da geração: um modelo que erra um cálculo tende
a errar também a conferência, e a literatura de \emph{self-consistency}
(Seção~\ref{sub:selfconsistency}) mostra que a repetição reduz, mas não elimina, erros
correlacionados. O SymPy oferece garantia de natureza diferente: para as classes tratáveis, o
erro do gabarito é \emph{detectado}, não apenas tornado menos provável.

\textbf{Por que orquestração em código puro, e não um \emph{framework} de agentes.} A
Seção~\ref{sub:orquestracao} apresentou LangGraph e CrewAI como candidatos. A implementação
optou por não adotar \emph{framework} algum: o fluxo tem quatro nós e uma política de decisão de
cinco regras, e expressá-lo diretamente em código torna a política legível por inteiro em uma
única função --- qualidade valiosa em um trabalho acadêmico cujo leitor é professor de
Matemática, não engenheiro de software. Um \emph{framework} se justificaria com fluxos maiores
ou paralelismo real entre avaliadores; a decisão é revisitada na discussão de trabalhos futuros.

\textbf{Por que JSON estruturado, e não texto livre.} Toda comunicação entre agentes usa objetos
validados por esquema. Texto livre exigiria que cada agente interpretasse a prosa do anterior
--- reintroduzindo, dentro do sistema, exatamente a ambiguidade que a arquitetura existe para
controlar. A validação por esquema também converte falhas de formato do LLM em erros detectáveis
e tratáveis, em vez de corrupção silenciosa de dados.
